\documentclass[xcolor=x11names,compress]{beamer}

%% General document %%%%%%%%%%%%%%%%%%%%%%%%%%%%%%%%%%
\usepackage{graphicx}
\graphicspath{{graphics/}}
\usepackage{tikz}
\usepackage{color,colortbl}
\usepackage{framed}
\usepackage{textcomp, setspace} %Needed for customization of ``listings''
% package
\usepackage[procnames]{listings} % to display code; don't forget [fragile]
% option after \begin{frame}
\input{inputs/rgb}
\definecolor{shadecolor}{rgb}{1,.9,.3}

\usepackage [autostyle]{csquotes}
\MakeOuterQuote{"}

\lstset{
    backgroundcolor=\color{shadecolor},
    tabsize=4,
    rulecolor=,
    language=python,
        basicstyle=\ttfamily\setstretch{1},
        upquote=true,
        aboveskip={1.5\baselineskip},
        columns=fixed,
        showstringspaces=false,
        extendedchars=true,
        breaklines=true,
        prebreak = \raisebox{0ex}[0ex][0ex]{\ensuremath{\hookleftarrow}},
        frame=single,
        showtabs=false,
        showspaces=false,
        showstringspaces=false,
        identifierstyle=\ttfamily,
        keywordstyle=\color[rgb]{0,0,1},
        commentstyle=\color[rgb]{0.133,0.545,0.133},
        stringstyle=\color[rgb]{0.627,0.126,0.941},
        numbers=left, 
        numberstyle=\tiny, 
        stepnumber=2, 
        numbersep=5pt
}

%%%%%%%%%%%%%%%%%%%%%%%%%%%%%%%%%%%%%%%%%%%%%%%%%%%%%%

%% Beamer Layout %%%%%%%%%%%%%%%%%%%%%%%%%%%%%%%%%%
\usetheme{Madrid}
\usecolortheme{seahorse}
\useoutertheme[subsection=false,shadow]{miniframes}
\useinnertheme{default}
% \usefonttheme{serif}
\usepackage{palatino}

\setbeamerfont{title like}{shape=\scshape}
\setbeamerfont{frametitle}{shape=\scshape, series = \bfseries}
\setbeamertemplate{frametitle}[default][center]
\setbeamertemplate{headline}{} %suppress headline (navigation pane)

\setbeamertemplate{footline}
{
  \leavevmode%
  \hbox{
  \begin{beamercolorbox}[wd=.1\paperwidth,ht=2.2ex,dp=1ex,center]{author in head/foot}
    \usebeamerfont{author in head/foot}\insertshortauthor
  \end{beamercolorbox}
  \begin{beamercolorbox}[wd=.8\paperwidth,ht=2.2ex,dp=1ex,center]{title in head/foot}%
    \usebeamerfont{title in head/foot}\insertshorttitle\hspace*{3em}
  \end{beamercolorbox}}
  \begin{beamercolorbox}[wd=.05\paperwidth,ht=2.2ex,dp=1ex,center]{}
     \insertframenumber{} / \inserttotalframenumber\hspace*{-3ex}
  \end{beamercolorbox}
}

\renewcommand{\(}{\begin{columns}}
\renewcommand{\)}{\end{columns}}
\newcommand{\<}[1]{\begin{column}{#1}}
\renewcommand{\>}{\end{column}}

\def\signed #1{{\leavevmode\unskip\nobreak\hfil\penalty50\hskip2em
  \hbox{}\nobreak\hfil(#1)%
  \parfillskip=0pt \finalhyphendemerits=0 \endgraf}}

\newsavebox\mybox
\newenvironment{aquote}[1]
  {\savebox\mybox{#1}\begin{quote}}
  {\signed{\usebox\mybox}\end{quote}}
%%%%%%%%%%%%%%%%%%%%%%%%%%%%%%%%%%%%%%%%%%%%%%%%%%%%
\title[Model Fitting]{Fitting Models to Data in Ecology and Evolution}
\author [Samraat]{Samraat Pawar}
\institute{{\it Department of Life Sciences (Silwood Park)}\\
  \vspace{12pt}
  \centering
  \includegraphics[height = .3in]{graphics/Imperial_Color1.pdf}
}  
%%%%%%%%%%%%%%%%%%%%%%%%%%%%%%%%%%%%%%%%%%%%%%%%%%%%

\begin{document}

%%%%%%%%%%%%%%%%%%%%%%%%%%%%%%%%%%%%%%%%%%%%%%%%%%%%%%
\begin{frame}[plain]
\titlepage
\end{frame}

% \title{Fitting Models to Data in Ecology and Evolution}
% % \vspace{12pt}
% % \subtitle{CMEE Masters}
% \author{
%     Samraat Pawar\\
%     \vspace{20pt}
%   \centering
%   \includegraphics[height = .3in]{Imperial_Color1.pdf}\\
% }
 

	
%%%%%%%%%%%%%%%%%%%%%%%%%%%%%%%%%%%%%%%%%%%%%%%%%%%%%%
\begin{frame}{Mechanistic vs. Phenomenological models}

	\begin{center}

		\it What does ``modelling data'' mean to you?
		
	\end{center}

 
 \end{frame}
 
 %%%%%%%%%%%%%%%%%%%%%%%%%%%%%%%%%%%%%%%%%%%%%%%%%%%%%%
\begin{frame}{Mechanistic vs. Phenomenological models}

	\begin{itemize}[<+->]\itemsep20pt

		\item {\it Mechanistic models} aim to explain the PROCESSES or
		MECHANISMS underlying PATTERNS or PHENOMENA in empirical data 
		\begin{itemize}
			\item These models have a THEORETICAL BASIS
		\end{itemize}
 
		\item {\it Empirical/Phenomenological models} establish the existence of
		STATISTICALLY SIGNIFICANT, NON-RANDOM PATTERNS or PHENOMENA in empirical
		data 
		\begin{itemize}
			\item They make no assumptions about the processes or mechanisms that generate the patterns
			\item That is, these models lack a THEORETICAL BASIS
		\end{itemize}
			
	  \end{itemize}
 
 \end{frame}
 
%%%%%%%%%%%%%%%%%%%%%%%%%%%%%%%%%%%%%%%%%%%%%%%%%%%%%%
\begin{frame}{Mechanistic vs. Phenomenological model fitting}

\begin{center}
		\includegraphics[width=.7\textwidth]{graphics/L-H_cycle.png}\\
		{\it \tiny source: \url{https://www.cds.caltech.edu/~murray/amwiki/images/8/8f/LHgraph.gif}}
\end{center}
\pause
\begin{itemize}[<+->] \itemsep6pt

	\item {\bf Mechanistic model}: \it The Lynx-Hare Cycle is driven by density-dependent population growth in hares

	\item {\bf Phenomenological model}: \it The Lynx and Hare Cycles have a significant asynchrony (period shift) of x years

\end{itemize}

 \end{frame}
 
 %%%%%%%%%%%%%%%%%%%%%%%%%%%%%%%%%%%%%%%%%%%%%%%%%%%%%%
\begin{frame}{Mechanistic vs. Phenomenological model fitting}

	\begin{itemize}[<+->]\itemsep10pt
		\item {\it It's not really one vs. the other}; Both types of models play a role in science (and Biology)
		\item Phenomenological model-fitting reveals patterns in data that generate HYPOTHESES 
 		\begin{itemize}
			\item These can be tested using further model fitting
			\item Example: {\it Whether} climatic temperature affects the Lynx-Hare cycle (using Generalized Linear Model-fitting)
		\end{itemize} 

		\item Mechanistic model-fitting {\it tries} to validate a mechanistic model that can explain the observed phenomenological pattern and generate MORE ACCURATE, MECHANISTIC HYPOTHESES
		\begin{itemize}
			\item Example: {\it How} climatic temperature {\it drives} the Lynx-Hare cycle
		\end{itemize} 
		\item \it Ultimately, successful, EMPIRICALLY-GROUNDED mechanistic models are the best path towards a THEORY in any scientific discipline (including ecology and evolution)  
	  \end{itemize}
 
 
 \end{frame}

  %%%%%%%%%%%%%%%%%%%%%%%%%%%%%%%%%%%%%%%%%%%%%%%%%%%%%%
\begin{frame}{Mechanistic vs. Phenomenological model fitting}
	 
  \begin{center}
	  \includegraphics[width=\textwidth]{Mechanisms.pdf}
  \end{center} 
 
 \end{frame}

 %%%%%%%%%%%%%%%%%%%%%%%%%%%%%%%%%%%%%%%%%%%%%%%%%%%%%%
\begin{frame}{Mechanistic models in Ecology and Evolution?}

\begin{itemize}[<+->] \itemsep12pt
	\item \it Do most ecological studies perform phenomenological or mechanistic modelling (or model-fitting)?
	
	\item The answer is mostly Phenomenological \pause --- {\it Why?} \pause
	\begin{itemize}
		\item Partly because we are still establishing the existence of GENERAL PATTERNS/PHENOMENA,
		\item  ... and partly because we are (or are forced to be) interested in FORECASTING rather than EXPLAINING. 
 
	\end{itemize}	  
	
	\item {\it So the big question is, can we FORECAST WITHOUT EXPLAINING?}
	
	\begin{itemize}
		\item For example, disease outbreaks: Do we really need to care about the underlying mechanisms if we can predict a future event using phenomenological modelling (e.g., Machine-learning of time series patterns)?
		
	\end{itemize}

\end{itemize}

\end{frame}

%%%%%%%%%%%%%%%%%%%%%%%%%%%%%%%%%%%%%%%%%%%%%%%%%%%%%%
\begin{frame}{What are Mechanisms?}

   \begin{itemize}[<+->] \itemsep4pt{}
			\item Somewhat subjective! 
			\item For example, the Ricker model can be thought of as 
			mechanistic:
			$$ 
				N_{t+1} = N_t e^{r\left(1-\frac{N_t}{k}\right)} 
			$$ 
				
			\item What is the mechanism? \pause --- Density dependence through scramble competition (Brannstrom \& Sumpter 2005)

			\item If the Ricker model and another model with contest 
			competition were compared with data --- some would call it mechanistic  modelling because one is trying to get at the underlying mechanism, scramble or contest competition
			
			\item But is this REALLY mechanistic? What are $r$ and $k$ really?
   \end{itemize}

\end{frame}

%%%%%%%%%%%%%%%%%%%%%%%%%%%%%%%%%%%%%%%%%%%%%%%%%%%%%%
\begin{frame}{Example of a fundamental mechanism: metabolic rate}

	\begin{itemize}
		\item Proponents of {\it Metabolic Theory of Ecology} argue that we have not progressed far enough towards mechanistic modelling because metabolism has been ignored
	\end{itemize}

	\begin{columns}[c]
  \column{.4\textwidth}
    \pause
      \includegraphics[width=\textwidth]{graphics/Photobacterium.pdf}
  \pause
  \column{.6\textwidth}
  \scriptsize
    $B = B_0 \boxed {e^{-\frac{E}{kT}}}f(T,T_{pk},E_D)$\\
    \vspace{10pt}
    \raggedright{$T$ = temperature (K)\\
     $k$ = Boltzmann constant (eV K$^{-1}$)}\\
     $E$ = Activation energy (eV)\\
     $T_{pk}$ = Temperature of peak performance\\
	 $E_D$ = Deactivation energy (eV)\\
	 
    {\tiny (J H van't Hoff 1884, S Arrhenius 1889)}
\end{columns}

\begin{itemize}\setlength{\itemindent}{0em}
  \pause
  \item Surely there is more to thermal responses?
    % \begin{itemize}\setlength{\itemindent}{-1em}
    %  \item Oxygen limitation
    %  \item Complexity of metabolic network
    %  \item Hormonal regulation
    % \end{itemize}
    % \pause
  \item \textit{What about alternative models?}\\

\end{itemize}
 
\end{frame}

%%%%%%%%%%%%%%%%%%%%%%%%%%%%%%%%%%%%%%%%%%%%%%%%%%%%%%
\begin{frame}{Modelling, and fitting models to data: What's the big idea?}

\begin{itemize} [<+->] \itemsep8pt

	\item {\it If possible}, use biological knowledge to construct models
	\item See if the models ``agree well'' with data
	\item Whichever model ``agrees best'' is most likely to have the right 
	mechanisms
	\item That's the one that's best for predictions (e.g. population 
	cycles), estimating rates (e.g. population or individual growth rates), etc
	\item Don't use models you already know have the wrong mechanisms just because they are popular!
	\item Phenomenological/statistical models often perform better than mechanistic ones. \pause {\it Why? \pause --- because they have lest restrictive assumptions}   

\end{itemize}
  
\end{frame}

%%%%%%%%%%%%%%%%%%%%%%%%%%%%%%%%%%%%%%%%%%%%%%%%%%%%%%

\begin{frame}{Models: How to build them?}

\begin{itemize}\itemsep10pt

	\item It's an art, takes practice 
	% (look at Levins' paper on the 
	% strategy of model building in biology)
	\item Build models one mechanism at a time --- in biology, it means 
	start at the right level of organization! 
	\item Always consider an alternative that is more parsimonious, even if 
	it is phenomenological (the thermal performance curves example: Sharpe-Schoolfield, Briere, or Polynomial?)! 

	\item For example, the Boltzmann-Arrhenius model is a good first try 
	describe and uncover mechanisms underlying individual level rates (e.g., vector fecundity or development rate)  

	\item The next step would be to include species interactions with  
	temperature dependence of individuals (or go in an evolutionary direction)

\end{itemize}   

\end{frame}

%%%%%%%%%%%%%%%%%%%%%%%%%%%%%%%%%%%%%%%%%%%%%%%%%%%%%%
\begin{frame}{Fitting models to data}

Multiple ways to do it:

\begin{itemize}\itemsep10pt

	\item Least Squares methods
	\begin{itemize}
		 \item Linear 
		 \item Non-linear
	\end{itemize}
	\item Likelihood-based methods 
	\begin{itemize}
	 \item Maximum Likelihood Estimation (MLE) 
	 \item Bayesian
	\end{itemize}
	\item Artificial intelligence and Machine learning
	\begin{itemize}
		\item  Focus in on maximizing ability to discover pattern and predict at the cost of mechanistic insights
	\end{itemize}

	% One-step forecasting (appropriate for discrete patterns/models), including time series data

\end{itemize}

\end{frame}

%%%%%%%%%%%%%%%%%%%%%%%%%%%%%%%%%%%%%%%%%%%%%%%%%%%%%%
\begin{frame}{Methods you can use}

\begin{itemize} \itemsep12pt
	
\item Least squares: along with Linear Model fitting, Non-linear Least Squares (NLLS) fitting is a particularly versatile and powerful approach because many mechanisms in biology and inherently non-linear

\item MLE/Bayesian methods: more robust if you are able to calculate the likelihood function analytically or numerically. 

\item AI/machine Learning: most versatile for large amounts of noisy data

\end{itemize}

\end{frame}


%%%%%%%%%%%%%%%%%%%%%%%%%%%%%%%%%%%%%%%%%%%%%%%%%%%%%%
\begin{frame}{Summary: Model selection is the key}

\begin{itemize}[<+->]\itemsep12pt
	\item Ideally, several competing (meaningful, not just null) hypotheses (mathematical models) should be fitted to data and compared using statistical theory 
	\item This is an advance over the traditional ``null hypothesis'' approach in Biology
	\item Necessary for developing the advancement of Biology from from an observational and axiomatic discipline to one with general theories.
	\item Necessary for understanding the mechanisms underlying biological patterns/phenomena
\end{itemize}
 
 \end{frame}

%%%%%%%%%%%%%%%%%%%%%%%%%%%%%%%%%%%%%%%%%%%%%%%%%%%%

\begin{frame}{Comparing and selecting models}

\begin{itemize}
	
	\item It's all about the ``Likelihood'' of a model: \\
	the set of parameter values of the	model ($\theta$) given outcomes ($x$),
	equals the probability of those observed outcomes given those parameter
	values, that is,

	$$ \mathcal{L}(\theta |x) = P(x | \theta)$$

	\item The easiest thing to do for you is to use information theory
	(including AIC and BIC) to compare models. 

	\item Both AIC and BIC use the {\it estimated likelihoods of a model}:\\
	
	AIC: $-2 \ln[\mathcal{L}(\theta |x)] + 2p$\\

	Small sample AIC (AICc): $-2 \ln[\mathcal{L}(\theta |x)] + 2p$\\
	
	BIC (Schwartz criterion): $-2 ln[\mathcal{L}(\theta |x) ] + p \ln(n)$
	
		(where $n = $ sample size, $p$ number of free parameters)
	\item The lower the AIC or BIC, the better. 

\end{itemize}

\end{frame}

%%%%%%%%%%%%%%%%%%%%%%%%%%%%%%%%%%%%%%%%%%%%%%%%%%%%%

\begin{frame}{Comparing and selecting models}

This is how you calculate AIC and BIC (using python syntax): 

\pause
\begin{itemize}\itemsep10pt
\small
	\item residuals = Observations - Predictions
	\item rss = sum(residuals ** 2) 
	\item Then, AIC is n * log((2 * pi) / n) + n + 2 + n * log(rss) + 2 * p \\
		({\it note n and p!})
	\item And BIC is n + n * log(2 * pi) + n * log(rss / n) + (log(n)) * (p + 1)

	\item That is, 	$ \mathcal{L}(\theta |x) = -\frac{n}{2/ln(RSS/n)}$
	\item For both AIC and BIC, If model {\bf A} has AIC lower by 2-3 or 
	more than model {\bf B}, it's better --- Differences of less than 2-3 
	don't really matter

\end{itemize}

    \pause
Also note that:
\begin{itemize}
\small
	\item R$^{2}$ = 1 - (rss$/$tss), where tss is total sum of squares: \\
tss = sum((Observations - mean(Predictions)) ** 2)\\
(a useful measure of goodness of fit)

\end{itemize}

\end{frame}

%%%%%%%%%%%%%%%%%%%%%%%%%%%%%%%%%%%%%%%%%%%%%%%%%%%%%

\begin{frame}{Comparing and selecting models: More stuff}

\begin{itemize} \itemsep12pt
	\item You can also calculate Akaike Weights, which is very useful/important
	when comparing $ > 2$ models. These weights can then be used to perform {\it
	model averaging}.

	\item Model selection using the Likelihood-Ratio test (LRT) is another
	option when you are comparing 2 models.
	
	\item Adjusted $R^2$ can be used to get rigorous ``idea'' about how
	alternative models are performing. 
	
	\item Very often, you will end up doing  model simplification, especially in
	{\it for linear least squares model fitting} --- starting with a complex
	model and then dropping terms till you have found a the most parsimonious
	version of the original model. There are functions in {\tt R} to do this (of
	course!).
 
\end{itemize}

\end{frame}

%%%%%%%%%%%%%%%%%%%%%%%%%%%%%%%%%%%%%%%%%%%%%%%%%%%%%%
\begin{frame}{Readings}

\begin{itemize}

\item Levins, R. (1966) The strategy of model building in population 
biology. Am. Sci. 54, 421--431.  

\item Johnson, J. B. \& Omland, K. S. (2004) Model selection in ecology 
and evolution. Trends Ecol. Evol. 19, 101--108. 

\item Bolker, B. M. et al.  (2013) Strategies for fitting nonlinear ecological models in R, AD Model Builder, and BUGS. Methods Ecol. Evol. 4, 501--512 .

\item Some illustrative examples of (non-linear) model fitting to 
ecological/evolutionary data 
\url{https://groups.nceas.ucsb.edu/non-linear-modeling/projects} 

\item Additional readings in the git repository
 
\end{itemize}
\end{frame}

\end{document}
